\documentclass[12pt]{article}
%\usepackage{Sweave}
\usepackage[margin=1in]{geometry}
%\usepackage[parfill]{parskip}
%\usepackage[hidelinks]{hyperref}
\usepackage[colorlinks = true,
            linkcolor = black,
            urlcolor  = blue,
            citecolor = black,
            anchorcolor = black]{hyperref}
\usepackage{rotating}
%\usepackage[backend=biber, style=apa, natbib=true]{biblatex}
\usepackage[citestyle=authoryear, style = apa, natbib=true, backend=bibtex, maxcitenames=2]{biblatex}
\addbibresource{references.bib}
\newtheorem{hyp}{H}
\newtheorem{hypo}{Hypothesis}
\usepackage[compact]{titlesec}
%\usepackage{titlesec}
\titleformat*{\section}{\bfseries\fontsize{14}{18}\selectfont}
\titleformat*{\subsection}{\bfseries\fontsize{12}{16}\selectfont}
%\usepackage{natbib}
\usepackage{csquotes}
\usepackage{xfrac}
\usepackage{mathtools}
\usepackage{booktabs}
\usepackage{lscape}
\usepackage{amsmath}
\usepackage[export]{adjustbox}
\usepackage{graphics}
\usepackage{array}
\usepackage[font=small,skip=1pt]{caption}
\usepackage{subcaption} 
\usepackage{ragged2e}
\usepackage{float}
\usepackage[section]{placeins}
\usepackage{hyperref}
\usepackage{multicol}
\usepackage{setspace}
\usepackage[hang,flushmargin]{footmisc}
\usepackage{afterpage}
\usepackage{soul}
\usepackage{xr}
\usepackage{booktabs}
\usepackage{multirow}
\usepackage{siunitx}
\usepackage{caption}
\usepackage{threeparttable}
\usepackage{makecell} 
\usepackage{tabularx}
\newcommand{\tabitem}{\textbullet~~}
\allowdisplaybreaks

% Set caption to be left-aligned
\captionsetup{justification=raggedright, singlelinecheck=false}

% Configure siunitx
\sisetup{
  group-separator = {,},
  group-minimum-digits = 4 % Start grouping from 4 digits
}

\doublespacing

\title{%
  The Signaling Gap: \\
  \large State-Owned Media and the Strategic Management of Sentiment during the Russian Invasion of Ukraine}
\author{Donald Moratz and Sarah Gerstein}
\date{\today}

%\title{The Signaling Gap: State-Owned Media and the Strategic Management of Sentiment during the Russian Invasion of Ukraine}

\begin{document}

\maketitle


\abstract{How do states strategically signal positions during international crises? This paper investigates the 2022 Russian invasion of Ukraine as a seminal shock to analyze the divergent responses of state-controlled and independent media in post-Soviet and satellite states. While independent media reflects unconstrained sentiment, we argue that state-owned media serves as a strategic signaling device. Utilizing a corpus of several hundred thousand articles, we employ a Large Language Model (GPT-4o) to isolate sentiment toward Russia from general war-related topical shifts—an approach that outperforms traditional BERT-based classifiers in capturing nuanced elite cues. Our findings reveal that while sentiment declined across our sample, the drop was significantly muted in state-owned outlets. Crucially, we find that states bound by Russian security agreements utilized media sentiment to signal strategic dissatisfaction even when formal diplomatic condemnation remained politically impossible. These results highlight the role of controlled media in communicating "true" intent during geopolitical upheaval and demonstrate the power of generative AI in identifying latent elite signaling within massive text corpora.}

\section*{Introduction}

How do constrained states signal genuine preferences when formal diplomatic channels are too costly to use? The Russian invasion of Ukraine in February 2022 serves as a seminal shock to the post-Cold War era. As NATO debated invasion preparations \citep{jakes_for_2024} and US political rhetoric questioned alliance commitments \citep{sanger_outburst_2024}, smaller states faced acute pressure to signal their strategic alignments without triggering costly retaliation from Moscow. One observable mechanism through which such signals can be transmitted is state media, which shapes public opinion through selective attribution and sentiment \citep{guriev_theory_2020, rozenas_how_2019}, yet the international relations literature has focused primarily on how these elite cues impact domestic audiences \citep{guisinger_mapping_2017}. Consequently, the divergence between state and private media sentiment, and the specific utility of media as an international signaling device, remains under-explored. We argue that state-owned media functions as a deniable informal signaling channel for states whose formal diplomatic constraints make official condemnation impossible, and that this dynamic is observable in the sentiment divergence between state and independent media during the 2022 invasion.

By using a generative AI model rather than traditional BERT-based classifiers, we isolate sentiment specifically toward Russia while controlling for topic. This ensures we capture shifts in attitude toward Russia itself, separate from the confounding influence of general war reporting. This approach allows us to identify how strategic elites responded to the crisis and clarifies the relationship between elites and foreign policy. We present three findings. First, the invasion caused a significant decline in media sentiment across all topics, separate from the shift toward war coverage. Second, we find that the causal effect of the invasion on sentiment differs significantly between independent and state-owned media, highlighting the strategic management of sentiment. Finally, we show that states with Russian security agreements expressed greater negative sentiment post-invasion than neutral countries, despite refusing to condemn the invasion publicly. This suggests these states utilize media to communicate genuine preferences when formal statements are politically impossible.

Signaling is a cornerstone of international relations literature. While theoretical work often centers on signaling credibility to rivals \citep{fearon_signaling_1997} or allies \citep{blankenship_trivial_2022}, research has largely focused on the domestic agenda-setting role of media \citep{mccombs_agenda-setting_2002}. Building on \citet{cohen_theatre_1987}, we contend that states strategically control sentiment toward Russia not only for domestic audiences but also to signal preferences to the international community. Such media signals reveal genuine preferences when formal diplomatic channels, such as UN votes, are constrained by high political costs. The invasion of Ukraine represents a crisis where states may publicly maintain one position while signaling a different intent through controlled media. This strategic behavior stands in contrast to independent media which remains unconstrained by state signaling imperatives.

Previously, data limitations hindered large-scale examination of these dynamics. However, advances in natural language processing provide new avenues for inquiry. We study how state-controlled media in states in Russia's near abroad responded to the invasion compared to independent sources. Utilizing a corpus of several hundred thousand articles from the period surrounding February 24, 2022, we applied a keyword-based approach to identify Russia-centric coverage. We then employed a generative AI model (GPT-4o) to conduct sentiment analysis, validating the model's performance against a human-coded gold standard. This approach allows us to causally identify how coverage changed around the discontinuity of the invasion for both state-owned and independent media, while controlling for topical shifts. The comparison across states with and without Russian security agreements is correlational, though the sharp discontinuity of the invasion provides a natural experiment that substantially reduces concerns about confounding within each group. %Our validation yielded a Krippendorff's alpha of 0.86, indicating high inter-coder reliability on an ordinal sentiment scale. 

The article proceeds as follows. We first develop our theoretical argument and derive three hypotheses about media divergence under alliance constraints. We then present the case, justifying our sample and identification strategy. The following sections describe our data and research design before presenting results. We conclude by discussing implications for how states signal preferences under geopolitical constraint.


\section*{Theory and Hypotheses}

Intuitively, the idea that state and independent media might react differently to a violation of an international border is not a novel concept. The two media sources might have very different mandates as well as funding sources. In addition, as a tool of the state, state-owned media is restricted by the states' decisions on how to make and signal alliance commitments. The international relations literature has grappled with several important questions related to alliances. Do states align with those most powerful and closest to them? Do states keep their friends closer, and their enemies closer? Upon deciding who to align with, how can small states, without powerful weapons and additional troops to send abroad, signal their commitment? The lack of concrete theory to help explain this variation in media response means that we must look to alliance, signaling, and media literature to help us begin to understand why we might see a difference in media sentiment. 

Weaker states often choose to balance in an effort to maintain their place in the system \citep{waltz_theory_2010}. When states do this, they are choosing the weaker side due to the perceived threat of the stronger state. States might also choose to bandwagon, that is join the stronger side \citep{waltz_theory_2010, walt_origins_2007}. States may choose to align with those who have the most aggregate offensive capabilities, who are located most closely, and who have the most aggressive perceived intentions. While bandwagoning might seem more attractive, Walt suggests that balancing is more common among states that matter and who have more capabilities \citep{walt_origins_2007}. 

While security concerns are likely to drive alliance preferences, states may also choose to join alliances that more closely adhere to their state views. Ideologically similar alliances are more common than alliances among countries of opposing ideological systems \citep{walt_origins_2007}. However, politicians might use ideological similarities to tout the strength of an alliance, though this level of closeness might be a facade, designed to inspire feelings of in-group favoritism. 

States may also enter into alliances designed to constrain. While some alliances are capability aggregation measures, meant to stand firm against external threats, alliances also may serve to tie squabbling partners together \citep{weitsman_dangerous_2004}. The alliance paradox suggests that allies are not only frequently unreliable, but also prone to confrontation. When states tether themselves by aligning with a powerful ally, they facilitate communication and mitigate potential hostilities, but this strategy constrains their partners. Tethering happens most often in relatively higher threat relationships; states that see themselves threatened by low levels of threat are more likely to hedge. States that are less threatened will "allow the hedging state to draw others into its sphere of influence as long as the course of action is not overly costly or too provocative to potential enemies" \citep{weitsman_dangerous_2004}. Small states located near powerful, threatening states might therefore choose to tether themselves to the regional pole, hoping to remain sovereign. 

Once states have decided on the best partners, how can they signal that they are committed? Signalling might involve words or deeds \citep{fuhrmann_signaling_2014}. Most political science literature focuses on how states might signal to rival states their level of capability and resolve \citep{fearon_signaling_1997, kertzer_resolve_2016}. Others have noted that signalling to allies is particularly difficult \citep{schelling_arms_1967}. Committed allies might make public statements or they might deploy troops or weapons to allied locations \citep{fuhrmann_signaling_2014, blankenship_trivial_2022}. 

By making public statements, countries could signal to their allies that they are committed to maintaining alliances. One of the most effective ways to broadcast these statements might be through the use of the media, including television, newspapers, and social media. This allows these statements to be public, helping to tie hands \citep{fearon_signaling_1997}. The news media is often used to set agendas and prime audiences as well as signal elite cues \citep{iyengar_news_2010, chong_framing_2007}. The news media can set agendas by identifying and publicizing issues for the mass public. The media can also prime by influencing the issues that the public pays attention to when they make judgements. Importantly, the media can help to transmit elite signals. There are two general types of media in most states: state-owned and private, or independent media. 

The distinction between state-owned and private media firms is unclear in the literature. Public choice theory suggests that government-owned media outlets “distort and manipulate information” to prop up the incumbent regimes, while privately owned and independent media sources offer alternative views with largely unbiased and accurate information \citep{djankov_who_2003}. State-owned media might also be more closely aligned with past authoritarian regimes \citep{moehler_whose_2011}. On the other hand, state-owned firms might provide more fair and balanced accounts, given their governmental affiliation and monetary backing \citep{djankov_who_2003}. Independent media, unhindered by government regulation may be more free to seek out the full story, or they may be backed by private firms that seek to provide a narrative more aligned with a desired end state. Privately owned media outlets might be more responsive to the public and open to opposing perspectives \citep{moehler_whose_2011}. Independent media in states that are more fully aligned with western values, with journalists interested in maintaining the freedom of the press that comes as a part of the international liberal order should be more likely to report more accurately. This will help them to maintain credibility as well as demonstrate their commitment to showing what is going on in the world. 

Within the sphere of state-owned media, there are two primary strategic threads discussed in the literature. The dominant strand has focused on the use of state-owned media in communication with domestic audiences. For example, \citet{mccombs_agenda-setting_2002} emphasizes the power of the news media to focus public attention and shape perspectives. Careful framing of issues within media can have a dramatic impact on public support for government actions \citep{gershkoff_shaping_2005}. Political parties wield significant control over their followers' opinions and can wield this influence to effect dramatic changes in support for policies \citep{slothuus_how_2021}. This is in conjunction with an emphasis within the broader discipline on domestic affairs \citep{tomz_public_2020, de_mesquita_testing_2004}. However, there has been a recent revitalization of research into the strategic use of media signalling to influence international audiences. 

Recent work has highlighted the extent to which savvy autocracies maintain power through the strategic manipulation of information--- management of state-owned media has become sophisticated, especially during crises \citep{guriev_spin_2022}. At the same time, there are important differences both across and within autocracies. Across autocracies, the underlying resource wealth of a state has been shown to have a profound impact on levels of media freedom \citep{egorov_why_2009}, which might help explain different incentives for state media manipulation and points to a need to control for the underlying media environment within a state. In particular, \citet{weeks_autocratic_2008} argues that different autocracies have systematic differences in their behavior in the international sphere. We argue that a result of these systematic differences is variation in their capacity for credible formal signaling. Those states with less capability for formal signaling may have an incentive to supplement their formal signals through informal channels that can convey preferences that formal diplomatic constraints make difficult to express directly. We argue that state-owned media is a plausible channel for this communication for two reasons. First, there is historical reliance by many countries in our sample on state-owned media for international signaling \citep{cohen_theatre_1987}. Second, the explicit rise in efforts by western countries to systematically study state media \citep{manor_digitalization_2019, usic_ic_2024} gives senders reason to believe that the signal will be received and correctly interpreted.

A natural concern is that observed divergence between state media and official positions simply reflects internal incoherence rather than strategic communication. Indeed, weak democracies and authoritarian regimes should not be viewed as monolithic with respect to their internal political structure. Authoritarians face multiple factions they must balance, a general public and a cadre of elite insiders \citep{svolik_politics_2012}. These elite insiders are not just an arm of the ruler, but a set of actors with their own preferences and survival calculations that require continuous bargaining and cooptation on behalf of the autocrat \citep{bueno_de_mesquita_logic_2003}. Dictators must manage elite bargaining \citep{gandhi_political_2008}, and we extend this argument to posit that state-owned media, as an outlet of the regime, faces pressure from both the political leader, but also from elites involved with the running of the outlet itself which may reflect different competing internal voices. Expressions of dissent even in the most constrained autocracies are consistent with different factional interests \citep{malesky_nodding_2010}. We argue that the signals sent by state-owned media should be considered important instruments of authoritarian communication and in particular, we should pay careful attention when the signals of state-owned media diverge from official positions. Whether state media diverges from formal positions because of top-down strategic choice or because of elite factional pressure, the divergence is analytically meaningful. While our data does not allow us to distinguish between these two mechanisms, our regression discontinuity design allows us to identify the causal effect of the invasion on state-media sentiment, independent of the diplomatic record. This means that divergence between the two is attributable to something other than common confounders, and we further develop our identification strategy in our methods section.

%While there is significant variation across the type of government in our dataset, we argue that in this particular crisis, the best predictor of their behavior is their alignment within the international sphere.


Governments have a strong incentive to capture media. \citet{gehlbach_government_2014} provides a workhorse theoretical model for strategic state-owned media. They argue that state-owned media is a strategic tool used by governments to advance a specific agenda and as a result, bias in state-owned media is generally higher than in independent media. While they argue that, with sufficient incentive to advance their agenda, the government is more likely to engage in capture of independent media, either directly or indirectly through paying off media, in a crisis like the invasion of Ukraine, either strategy is unlikely to be immediately actionable in the short-term. Building upon this logic, we argue that we should expect to see a strong divergence between state-owned media coverage and independent media coverage initially around the invasion as state-owned media responds strategically and before governments can begin to influence independent media coverage. 

Importantly, state-owned media is effective in shaping audience opinions in weak democracies. \citet{enikolopov_media_2011} provides empirical evidence of the impact of independent media coverage. They show that exposure to independent media dramatically reduces vote-share for the ruling party in the 1999 Russian election. This highlights the extent to which state-owned media framing can matter in the absence of exposure to independent media--- it serves not just as a tool for mobilization of existing preferences, but can serve to shape and change attitudes. Importantly, they argue that media plays a powerful role in persuasion in countries with weak democratic institutions, which characterizes most of the countries in our sample. 

There is evidence that state-owned media has been used as a geopolitical instrument in the region. \citet{peisakhin_electoral_2018} use variation in exposure to Russian state TV to examine support for the Pro-Russian candidates in Ukraine's 2014 elections. They find that exposure increases support for Pro-Russian candidates, but primarily amongst citizens already inclined to support those candidates. This has two interesting implications for our paper; first, it provides additional evidence of the importance of the strategic use of state-owned media in the region, and second, it highlights that state-owned media can serve as an important geopolitical instrument.

Within the region, \citet{oates_television_2006} argues that state-owned media serves most directly as a tool for elites to engage in agenda setting. She highlights how effective it has been in managing political discourse and short-circuiting efforts to promote democratic institutions. This suggests that there are benefits to state-owned media strategically manipulating sentiment around an event like the invasion of Ukraine. This is particularly important within the context of the contemporary media ecosystem with readily available online spheres. Careful use of disinformation can provide advantages to repressive regimes seeking to manipulate audience opinions \citep{oates_russian_2016}.

The relationship between media and foreign policy has been examined through several lenses. The foundational text on media and foreign policy, \citet{robinson_cnn_2002} explored the idea of, and ultimately found mixed evidence for, a "CNN effect", whereby media coverage of crises forces government to intervene. Instead, he found that media influence is only likely when policy is uncertain and coverage is both critical of government and empathetic of those suffering during the crisis. Once policy is settled, he argues that media largely follows elite cues. We reframe this; we argue that rather than simply following elite cues or driving foreign policy, states with access to state-owned media strategically weaponize media as foreign policy.

\citet{entman_projections_2004} introduced a cascade model of media, with frames originating with the executive and filtering down through elites, media, and finally to the public, with friction at each level within a democracy. We build on his framework of media framing as a modern foreign policy instrument by extending his work to autocracies and weak democracies, where we argue that there is less friction at each level, especially with respect to state-owned media. This is driven by a combination of editorial processes, which explicitly work to build on official state positions, and journalistic incentives, where journalists face potentially harsh sanctioning for challenging the state position.

\citet{cohen_theatre_1987} emphasizes the role of diplomatic signaling through the media. He argues that these media signals are are ambiguous by design, allowing deniability, but are well understood by the intended recipients (other states)---- a dynamic that has only intensified as foreign governments now systematically monitor state media output \citep{manor_digitalization_2019, bjola_revisiting_2018, usic_ic_2024}. Indeed, during the Cold War, western analysts meticulously studied the state-owned media of the USSR and its satellites to understand the Soviet foreign policy. Soviet officials strategically used this fact to signal to the West in ways that they could not do through formal channels \citep{cohen_theatre_1987}. This work serves as a precursor to \citet{miskimmon_strategic_2013}, which introduces the concept of strategic narratives, which states leverage through media as a core instrument of foreign policy and serve to shape how states project their narrative to the world. Their framework establishes how actors form and project narratives at three levels; system, identity and issue. 

The need for informal signaling channels is particularly acute for autocracies. Formal signaling theory establishes that credible diplomatic signals require that the signaling state incur costs that make bluffing irrational, either by tying hands through audience costs or by sinking material costs \citep{fearon_signaling_1997}. Democratic governments can generate audience costs through open political competition, where visible domestic dissent or support provides external observers with reliable information about a government's true resolve \citep{schultz_democracy_2001}. Autocracies and weak democracies, which suppress dissent and coerce apparent consensus, face structural credibility deficits in formal channels: their official diplomatic statements are less informative precisely because external observers know they carry little or no domestic political cost. State media, in this context, becomes a more credible--- albeit still deniable--- channel because foreign governments monitor it as a less-controlled reflection of elite preferences \citep{bjola_revisiting_2018}. The digitalization of public diplomacy has formalized this dynamic. State media is now systematically tracked by foreign intelligence and diplomatic services as informative about the actual policy preferences of the sending government \citep{manor_digitalization_2019, usic_ic_2024}.

We extend the framework proposed by \citet{miskimmon_strategic_2013} by considering the constrained nature of international signaling due to a combination of trade and security integrations. We posit that strategic narrative projection operates at both formal levels (UN votes, official diplomatic statements) and more informal levels (state-owned media). We argue that the invasion of Ukraine triggers crises at all three levels of the framework in \citet{miskimmon_strategic_2013}. While the invasion serves foremost as an issue narrative contest, the implications of a major land war in Europe, an international order shaking crisis, triggers both system and identity pressures for post-Soviet states. The invasion is a direct challenge to the established post-Cold War system, challenging norms around the use of force and wars of territorial expansion. At the identity level, it challenged state positions--- its values, its alliances and its place in the world. For post-Soviet states that found themselves caught between Russian security agreements and Western integration, this introduced enormous identity issues. That this crisis provoked challenges around all three levels of strategic narratives introduced a tension that states needed to manage; while many states may have wished to project support for a sovereign state facing invasion, the system and identity narratives within a state may serve as a constraint to how far it can go. This serves as the crux of our third hypothesis--- states with Russian security agreements were forced to rely on issue-level media signals to express preferences that identity- and system-level constraints made impossible to communicate through formal diplomatic channels.

We build on this literature to examine strategic sentiment manipulation in state-owned media. Our theoretical argument assumes that the invasion of Ukraine represented a genuine negative shock to media sentiment across the region — without this baseline, the divergence we identify in H2 and H3 would be uninterpretable. We first establish this empirically:

\textbf{Hypothesis 1:} The invasion of Ukraine produced a significant decline in media sentiment toward Russia across both state-owned and independent outlets in our sample.

With this baseline established, we turn to the divergence between state-owned and independent media. State-owned media in our sample operates under political constraints that independent media does not face. We therefore expect that state-owned media will manage coverage more carefully than independent media, producing a smaller decline in sentiment toward Russia after the invasion.

\textbf{Hypothesis 2:} State-owned media will report negative sentiment towards Russia after the invasion to a lesser extent than independent media. 

A straightforward domestic audience management account would predict that states with formal security agreements with Russia suppress critical coverage of the invasion — protecting the alliance relationship means keeping the domestic public from reading negative coverage of Russia. Our signaling argument predicts the opposite: precisely because formal diplomatic condemnation is too costly, these states use state-owned media as an informal channel to express dissatisfaction to external audiences, producing more negative coverage than neutral states whose formal channels remain open. That the domestic account and the signaling account generate opposite predictions for H3 provides some leverage for distinguishing between them. If domestic audience management were the primary driver, we would further expect coverage patterns to track the size of ethnic Russian populations across our sample — since Russia explicitly justified the invasion as protection of ethnic Russians abroad, regimes with larger ethnic Russian minorities would face the strongest domestic incentive to suppress critical coverage. We return to the opposite-predictions test and the ethnic Russian robustness check in our discussion.

\textbf{Hypothesis 3:} State-owned media in states with formal security agreements with Russia will respond more negatively to the invasion than state-owned media in neutral states.

\section*{The Case}

On 24 February 2022, Russia launched a full-scale invasion of Ukraine, initiating the largest land war in Europe since World War II. The invasion was a direct violation of the UN Charter's prohibition on the use of force against the territorial integrity of a sovereign state \citep{nations_charter_1945}--- a norm that Russia, as a permanent member of the Security Council, had formally committed to uphold. The international community largely condemned the attack. The Council of Europe suspended Russian participation, and the UN Security Council took up a draft resolution condemning the invasion, which Russia vetoed \citep{walker_conflict_2024}. For most states, the path to a public response was straightforward. For the states in our sample, it was not.

The twelve states we study occupy very different positions relative to Russia, and this variation is central to our empirical design. We organize them into three broad categories based on their security and economic relationships with Russia at the time of the invasion. The first category comprises states with formal Russian security agreements: Armenia, Belarus, and Kazakhstan are all members of the Collective Security Treaty Organization (CSTO), Russia's primary multilateral security alliance in the post-Soviet space. Azerbaijan, while not a CSTO member, signed a Declaration on Allied Interaction with Russia in February 2022, one day prior to the invasion \citep{huseynov_azerbaijan_2022}, placing it in a similar position of formal alignment. These states faced the greatest constraints on formal condemnation; openly criticizing a security patron carries real costs, and CSTO membership in particular creates institutional pressure toward solidarity. At the same time, all four states also maintained substantial trade relationships with Western nations, giving them a real economic interest in condemning the invasion even where formal security commitments constrained them from doing so\footnote{The level of trade reliance expressed as a percent of total trade (imports plus exports) originating from the west (Europe and North America) for each country in 2021 was Belarus: 22\%, Armenia: 34\%, Kazakhstan: 46\%, and Azerbaijan: 72\% \citep{world_bank_world_2021}}.

The second category comprises states we treat as broadly neutral: Georgia, Moldova, Kosovo, and Serbia. Each of these states has maintained distance from both Russian security structures and full Western integration, though for different reasons and with different relationships to Russia's sphere of influence. Georgia and Moldova have both faced Russian territorial incursions in recent times\citep{cheterian_august_2009, potter_unrecognized_2022}, while Serbia and Kosovo's distance from both Russia and western integration reflects different dynamics entirely. Serbia has traditional ties to Russia, driven in part by its historical ties to Moscow that has created tension between its EU candidacy and its stance on the war \citep{radeljic_inefficiency_2023}. Kosovo, meanwhile, is dependent on western recognition for its own sovereignty claim, though it sits outside formal security structures on both sides \citep{newman_foreign_2016}. While our neutral states reflect a broad mix of relationships both with the West and with Russia, this mix allows them to function as an appropriate baseline of comparison for states with explicit security agreements.

The third category comprises states aligned with Western institutions: Albania and Hungary are both full NATO members, and Ukraine, while not a NATO member at the time of the invasion, had formally declared NATO membership a strategic objective. Turkey presents a more complex case: formally a NATO member since 1952, Ankara declined to join Western sanctions against Russia and instead positioned itself as a mediator between the two parties, reflecting the increasingly independent foreign policy orientation it had pursued since normalizing relations with Moscow after 2016 \citep{mankoff_turkeys_2022}. We treat Turkey as NATO-aligned for the purposes of our sample construction, while acknowledging that its behavior during the invasion did not fully mirror that of other alliance members. Ukraine is included in our dataset but excluded from the main analysis given its direct role as the invaded party. We describe our media sources for each country in the data section.

The constraints facing Russian-aligned states in February 2022 did not emerge overnight. Russia's willingness to use force in its near abroad had been established clearly by 2014, when Russia annexed Crimea following political upheaval in Ukraine and backed separatist forces in the Donbas. The Minsk Agreements, brokered by Western nations, called for a ceasefire and reaffirmation of Ukrainian sovereignty, but violence continued to escalate through the intervening years \citep{walker_conflict_2023}. By November 2021, Russia had mobilized troops along the Ukrainian border, and intelligence from both the U.S. and U.K. was indicating that a full-scale invasion was imminent \citep{walker_conflict_2023}. Russia simultaneously demanded that NATO commit to admitting no new members--- explicitly naming Ukraine and Georgia--- signaling that states in its near abroad would face consequences for Western alignment \citep{al_jazeera_us_2022}. For CSTO members and Russian-aligned states, this context made formal diplomatic condemnation of the invasion not just politically costly but institutionally constrained. State-owned media in the post-Soviet space has historically functioned as a tool for managing exactly these kinds of politically sensitive moments, projecting carefully calibrated signals to both domestic and international audiences while preserving official deniability \citep{cohen_theatre_1987, oates_television_2006, guriev_spin_2022}.

February 24 nonetheless represents a discrete and dramatic rupture in the information environment. While the troop buildup had proceeded for months, the full-scale invasion--- with missile strikes on Kyiv and ground forces crossing the border simultaneously--- was of a different order than the slow escalation that preceded it. Russia had conducted large-scale military exercises with Belarus near the Ukrainian border on February 10, with western intelligence reporting up to 140,000 troops positioned at the border \citep{walker_conflict_2024}. Even so, the decision to launch a comprehensive invasion rather than a more limited operation represented a threshold that most observers, including many within the states in our sample, had not fully anticipated. This discontinuity is the basis for our identification strategy: we treat February 24 as a sharp break in the media environment and use the variation in coverage before and after it to identify the causal effect of the invasion on state media sentiment. We develop this design in detail in the methods section.

\section*{The Data}

To study how the invasion changes sentiment towards Russia in state-sponsored media versus independent media, we utilized a corpus of news articles collected as part of the \href{https://web.sas.upenn.edu/mlp-devlab/}{Machine Learning for Peace Project}. These articles were collected via an extensive scraping process that targets both international and domestic news sources around the world. As part of the collection process, we traced ownership of each newspaper to identify those which are either directly state-owned or are indirectly captured by the state. Both independent and state-owned newspapers were carefully scraped via custom parsers in their original language and then human coders checked each newspaper for errors in the collection process. We then utilized a system of \href{https://huggingface.co/}{Hugging Face Transformers} to translate articles into English. These translations have undergone human validation to ensure accuracy and usability. Once articles have been translated, we utilized a carefully curated keyword list of over 500 words associated with Russia, including prominent politicians and companies, to identify articles related to Russia.\footnote{For more information on the process of article collection, please see \href{https://web.sas.upenn.edu/mlp-devlab/technical-details-research/}{the MLP Pipeline Report.}}

We collected articles from the 6 months on either side of the invasion, spanning from August 1$^{\text{st}}$, 2021 till August 30$^{\text{th}}$, 2022. Our curration process yielded 225,000 articles from 12 countries\footnote{Albania, Azerbaijan, Armenia, Belarus, Hungary, Kazakhstan, Kosovo, Moldova, Serbia, Turkey, Ukraine, and Georgia} and 70 sources, including at least one state-owned media source for each country. We chose these 12 countries according to two primary criteria. First, we chose a selection of former USSR and USSR-aligned states. Second, to ensure that Hungary, as a current NATO ally, was not unique, we added Turkey and Albania. Finally, we included Serbia as a non-NATO, non-former-USSR country, but with a government that has been decidedly pro-Russian in recent history to expand the analysis. For our main specifications in this article, we omit Ukraine to highlight the behavior in countries not directly affected by the invasion, but the inclusion of Ukraine does not meaningfully change the results.\footnote{To be included in future appendix.} After omitting Ukraine, we were left with 174,000 articles from 62 sources. 

To conduct our analysis, we utilized a zero-shot GPT 4.0 Turbo model to code articles across three dimensions; article topic (1 of 5), whether the article features Russia as a main actor (binary), and the sentiment of the article towards Russia (ternary). This coding was applied to the article title as well as the first two sentences.\footnote{Unlike the GPT classifier, the human codings utilized both the title as well as the whole article. There were not significant performance gains from applying the GPT classifier to the whole article, and thus the classifier was applied to a truncated portion of the article to save both computational and monetary costs.} Russia was considered to encompass both the broader state, as well as features of the state apparatus including the Russian military and ministers of the Russian government as well as the head of the state, Vladimir Putin. The topics were classified as one of ``Political,'' ``Economic,'' ``Military,'' ``Culture,'' or ``Other.'' Sentiment was allowed to take the values of -1 (negative), 0 (neutral) or 1 (positive). To ensure accuracy, we compared the model's performance against a test set of 100 randomly selected human-coded articles. The human coding process was adversarial double-coded by the coauthors who then debated the codings to settle on a final coding for each dimension for each article. The overall agreement rate for initial human-codings was 79\% with respect to sentiment, 91\% with respect to Russian-Related, and 93\% for topic. The final results from the GPT model with respect to sentiment are presented in Figure 1. Overall, the F1 score for the sentiment classifier is .92, for the Russian involvement classifier is .83, and for the Topic classifier is .80, with the confusion matrices reported in the appendix. The classifier raises two areas of concern. First, it is overly likely to classify political articles as other compared to human-coders. Second, it is overly likely to classify articles as not having Russian involvement compared to human-coders. Neither are likely to bias the findings reported in this article as any bias resulting from either issue would be towards a null result.

\begin{figure}[H]
    \centering
    \includegraphics[width=14cm]{Figures/figure_1.PNG}
\end{figure}

\newpage

\printbibliography

%\bibliographystyle{apalike}
%\bibliography{references.bib}

\end{document}